\section{Norm geometry and finite realization}\label{sec:normgeometry}

We use the complete space of all norms to construct a minimum and return
to split norms only after establishing finite regularity. This section
records the spectral convexity, finite realization and projection results
needed for that argument.

\subsection{Relative spectra and the completion}
For a finite-dimensional space $W$ over a complete non-Archimedean field,
write $\Ndiag(W)$ for the split norms and $\Nall(W)$ for all norms.
The Goldman--Iwahori distance is
\[
 d_\infty(h,k)=\sup_{z\ne0}\left|\log\frac{k(z)}{h(z)}\right|.
\]
If $h,k$ are split, choose a common orthogonal basis and arrange the
logarithms $\log(k(e_i)/h(e_i))$ in decreasing order:
\[
 \sigma(h,k)=(\sigma_1(h,k),\ldots,\sigma_n(h,k)).
\]
Put $S_j(h,k)=\sum_{i=1}^j\sigma_i(h,k)$ and $S_0=0$.
Equivalently, $S_j$ is the logarithm of the operator norm of the identity
on the $j$th exterior power, with its induced determinant norms.
This characterization extends $S_j$ continuously to all norms. Indeed,
if $d_\infty(h,h')\le\eps$ and $d_\infty(k,k')\le\eta$, comparison of
exterior norms gives
\begin{equation}\label{eq:spectralcontinuity}
 |S_j(h,k)-S_j(h',k')|\le j(\eps+\eta).
\end{equation}
We define the spectral coordinates on all norms by $\sigma_j=S_j-S_{j-1}$.

Fix a nonzero volume vector $\xi\in\bigwedge^nW$ and let $d(h)$ be the
logarithm of its determinant norm. For a basis matrix $P$ relative to a
fixed frame with volume $\xi$, and weights $w_i$, this is
\begin{equation}\label{eq:detnorm}
 d(h)=\sum_iw_i+\nu(\det P).
\end{equation}
Thus
\begin{equation}\label{eq:detspectrum}
 S_n(h,k)=d(k)-d(h),\qquad
 d_2(h,k)^2=\sum_i\sigma_i(h,k)^2,
 \qquad d_\infty\le d_2\le\sqrt n\,d_\infty.
\end{equation}
Changing $\xi$ adds the same constant to all determinant logarithms.
Only differences with a common volume convention will occur.

The canonical geodesic between two split norms interpolates their
logarithmic weights in a common orthogonal basis. Split norms are dense
in $\Nall(W)$; the geodesics and the spectral metric extend to the
completion. The space $(\Nall(W),d_2)$ is complete and CAT(0), and these
are its unique geodesics \cite[Theorem~1.19 and Corollary~3.3]{BE21}.
We always mean these geodesics, even when discussing a different
spectral metric. Determinant logarithms are affine along them.
The existence of a common basis for two split norms, as well as the
restriction and quotient properties of split norms, is given in
\cite[Lemma~1.13 and Proposition~1.14]{BE21}.

For $W=V$, use the product metric
$d((h_A,h_C),(k_A,k_C))^2=d_2(h_A,k_A)^2+d_2(h_C,k_C)^2$ when applying
CAT(0) arguments. Suprema in time will be taken with the equivalent
maximum of the two vertexwise distances. Let
\[
 \Dord=\{(h_A,h_C)\in\Nall(V)^2:h_A\le h_C\}.
\]

\begin{lemma}\label{lem:spectralconvexity}
Let $h_\theta,k_\theta$ be canonical geodesics. Then
\begin{equation}\label{eq:jointmajorization}
 \sigma(h_\theta,k_\theta)
 \prec(1-\theta)\sigma(h_0,k_0)+\theta\sigma(h_1,k_1),
 \qquad 0\le\theta\le1,
\end{equation}
where $\prec$ denotes majorization with equality of total sums.
Consequently, for every lower-semicontinuous convex function
$\phi:\RR\to\RR\cup\{+\infty\}$,
\begin{equation}\label{eq:costconvexity}
 (h,k)\longmapsto\sum_i\phi(\sigma_i(h,k))
\end{equation}
is jointly geodesically convex and lower semicontinuous.
The set $\Dord$ is closed and geodesically convex. Moreover,
\begin{equation}\label{eq:trianglemajorization}
 \sigma(h,l)\prec\sigma(h,k)+\sigma(k,l).
\end{equation}
\end{lemma}

\begin{proof}
First suppose the endpoints are split. Fix one basis for the geodesic
$h_\theta$ and another for $k_\theta$. The logarithm of the norm of an
exterior power of the identity is the maximum of finitely many affine
functions of the two sets of weights, indexed by its nonzero minors.
Thus every $S_j(h_\theta,k_\theta)$ is convex; this is the minor proof
of \cite[Lemma~6.16]{HKKP26}. For $j=n$, the determinant identity gives
affineness. The weighted sum on the right of
\eqref{eq:jointmajorization} is decreasing, so these are exactly the
majorization inequalities. The convex-hull characterization of
majorization and scalar Jensen's inequality give
\eqref{eq:costconvexity}. Neither monotonicity nor evenness of $\phi$
is needed. If $\phi$ is finite only on an interval, finite endpoint
costs and majorization keep the intermediate spectra in that interval;
the extended-real inequality follows as well.

Approximate the endpoints by split norms. Continuity of geodesics and
\eqref{eq:spectralcontinuity} extend majorization to all norms. The
convexity conclusion follows from the same majorization inequality,
and lower semicontinuity follows from that of $\phi$.
The order condition is equivalent to $\sigma_1(h_C,h_A)\le0$, so its
sublevel set is closed and convex. Finally, submultiplicativity of
exterior operator norms gives
$S_j(h,l)\le S_j(h,k)+S_j(k,l)$; equality for $j=n$ follows from
\eqref{eq:detspectrum}. This proves \eqref{eq:trianglemajorization}.
\end{proof}

We also use the following quantitative CAT(0) inequality. For arbitrary
$x,x',y,y'$ and their respective midpoints $m_x,m_y$,
\begin{equation}\label{eq:fourpoint}
 d(m_x,m_y)^2\le
 \frac{d(x,y)^2+d(x',y')^2}{2}
 -\frac{\bigl(d(x,x')-d(y,y')\bigr)^2}{4}.
\end{equation}
It is a form of the energy convexity inequality of
Korevaar--Schoen \cite[Proposition~1.3, equation~(1.3.2)]{KS97}. To see the coefficient,
write $b=d(x,x')$, $c=d(y,y')$. Two applications of the CAT(0)
midpoint inequality bound $4d(m_x,m_y)^2$ by
\[
 d(x,y)^2+d(x,y')^2+d(x',y)^2+d(x',y')^2-b^2-c^2.
\]
The CAT(0) quadrilateral inequality bounds the two cross terms by
$d(x,y)^2+d(x',y')^2+2bc$, giving \eqref{eq:fourpoint}.
No local compactness is required.

\subsection{Scalar extension and finite realization}
The field $K$ is algebraically closed
\cite[Appendix~A, Lemma~A.1]{FOOO10}. Let
\[
 L=\CC((T^\RR))
\]
now denote the Hahn field, whose supports are arbitrary well-ordered
subsets of $\RR$. It is algebraically closed and spherically complete
\cite[Lemma~4, Theorem~1 and Corollary~4]{Poonen93}. The natural inclusion $K\subset L$ preserves the
value group and residue field. Model completeness of algebraically
closed valued fields makes this inclusion elementary. In particular,
a finite existential condition over $K$ holding in $L$ holds in $K$;
see \cite[Definitions~2.1--2.2 and Theorem~3.10]{Yin10} for the valued-field
language and quantifier-elimination statement.

For $h\in\Nall(W)$, its canonical extension to $W_L$ is
\begin{equation}\label{eq:scalarextension}
 h_L(z)=\inf_{z=\sum_i b_i z_i}\max_i |b_i|h(z_i),
 \qquad b_i\in L,\quad z_i\in W.
\end{equation}
This is a norm, restricts to $h$, is isometric for $d_\infty$, and
preserves orthogonal bases and their weights
\cite[Definition~1.24 and Proposition~1.25]{BE21}.
Every finite-dimensional $L$-norm is split \cite[Lemma~1.12]{BE21}. Approximating a pair of
$K$-norms by split norms and using a common basis for the approximants
shows, by \eqref{eq:spectralcontinuity}, that canonical extension
preserves their entire relative spectrum. The same argument applies
to an invertible $K$-linear arrow by pullback. Consequently extension
is an isometry for $d_2$ and preserves canonical geodesics and order.

\begin{lemma}[Finite realization]\label{lem:finiterealization}
Let a finite directed graph have an $n$-dimensional $K$-space $V_v$ at
each vertex and a fixed $K$-linear isomorphism $T_e$ on each edge.
Give each $V_v$ an arbitrary norm $h_v$, and specify a subset of
vertices at which the norms are split. There are split $K$-norms $k_v$
which agree exactly with the prescribed norms at those vertices and
preserve the entire logarithmic singular-value list of every $T_e$.
\end{lemma}

\begin{proof}
Extend the diagram canonically to $L$, preserving every spectrum.
Choose orthonormal $L$-basis matrices $P_v$, using fixed orthonormal
$K$-basis matrices at the specified vertices. Orthonormal bases exist
because the norms are split over $L$ and both value groups are $\RR$.
Introduce matrices $Q_v$ with equations $P_vQ_v=Q_vP_v=I$.
The arrow in these bases is $M_e=Q_{t(e)}T_eP_{s(e)}$.

For each $1\le j\le n$ set
$c_{e,j}=-\sum_{i=1}^j\sigma_i(T_e)$. Preserving the first $j$
singular values in sum means precisely that every $j$-minor of $M_e$
has valuation at least $c_{e,j}$ and one such minor has valuation
exactly $c_{e,j}$. These are finite valued-field conditions with the
parameters $T^{c_{e,j}}\in K$: every normalized minor belongs to
$\OO$, and at least one is a unit. The boundary matrices and their
inverses are fixed, not merely the lengths of their columns.

The chosen $L$-matrices solve this finite formula. Elementarity gives
$K$-matrices solving the same formula. Declare their columns
orthonormal. The resulting norms have the required full spectra by
the exterior-minor formula, and the fixed boundary matrices give
exactly the prescribed boundary norms.
\end{proof}

Only finitely many real exponents are parameters in this argument.
It does not quantify over norm functions or over a continuum of time
values. Nor does it assert proximity between the newly realized norm
and the original norm at an unspecified vertex. These distinctions
will matter when passing to a limit of polygons.

\subsection{Projections associated with a decomposition}
Fix a direct sum $V=V_1\oplus\cdots\oplus V_m$. Let $\mathcal F$ be
the set of norms for which this decomposition is orthogonal:
\begin{equation}\label{eq:orthogonalblocks}
 h(z_1+\cdots+z_m)=\max_i h(z_i),\qquad z_i\in V_i.
\end{equation}
This permits arbitrary norms on the summands and does not require each
summand norm to be split. If the summands are lines, $\mathcal F$ is
the corresponding complete apartment.

\begin{lemma}\label{lem:blockprojection}
The nearest-point projection $P:\Nall(V)\to\mathcal F$ preserves
norm order and determinant logarithms, sends split norms to split norms,
and satisfies
\begin{equation}\label{eq:projectioncost}
 \sum_i\phi\bigl(\sigma_i(Ph,Pk)\bigr)
 \le\sum_i\phi\bigl(\sigma_i(h,k)\bigr)
\end{equation}
for every lower-semicontinuous convex scalar function $\phi$.
\end{lemma}

\begin{proof}
Let $G$ be generated by the signs $s_i$ which act as $-1$ on $V_i$
and as $1$ on the other summands. Since $|2|=1$, a $G$-invariant
norm makes the coordinate projection $(I-s_i)/2$ non-expansive.
Together with the ultrametric triangle inequality this proves
\eqref{eq:orthogonalblocks}. The converse is immediate. Thus
$\mathcal F$ is the common fixed set of a finite group of isometries;
it is nonempty, closed and convex.

The unique barycenter of the orbit $\{g_*h:g\in G\}$ belongs to
$\mathcal F$ by invariance. On $\mathcal F$ its squared-distance
objective is $|G|d_2(z,h)^2$, so it is precisely $Ph$.
A finite squared-distance barycenter of points in a closed convex set
lies in that set. Indeed, projection to the set decreases every squared
distance by at least the squared distance of the barycenter to its
projection. This would strictly decrease the objective unless the
barycenter already belonged to the set.

If $h_A\le h_C$, the orbit pairs belong to the closed convex order
set $\Dord$. Their product barycenter is $(Ph_A,Ph_C)$, so order is
preserved. The determinant level set through $h$ is closed and convex
and contains its orbit, because $|\det g|=1$. Hence $d(Ph)=d(h)$.
By Lemma~\ref{lem:spectralconvexity}, every finite sublevel set of
\eqref{eq:costconvexity} is closed and convex. The orbit pairs
$(g_*h,g_*k)$ have the same relative spectra and belong to this set.
Taking their barycenter proves \eqref{eq:projectioncost}. The case of
infinite right-hand side is automatic.

Finally suppose $h$ is split. Its orbit norms are split. Apply
Lemma~\ref{lem:finiterealization} to the unknown norm $Ph$, the fixed
orbit norms, and the identity arrows joining it to each orbit norm.
There is a split norm with the same distances to all orbit norms and
therefore the same barycenter objective. Uniqueness identifies it
with $Ph$.
\end{proof}

The contraction concerns the full convex spectral cost, rather than
only $d_2$. We shall also need its quantitative part. Extend
$\ell(v)=-\log\cos v$ by $+\infty$ for $|v|\ge\pi/2$.
For $\Delta,W>0$ set
\[
 \mathcal K_{\Delta,W}(h,k)
       =\frac W2\sum_i\ell\bigl(\sigma_i(h,k)/\Delta\bigr).
\]

\begin{lemma}\label{lem:projectionloss}
For a finite-cost pair $h,k$, the projection of
Lemma~\ref{lem:blockprojection} satisfies
\begin{equation}\label{eq:projectionloss}
 \mathcal K_{\Delta,W}(h,k)-\mathcal K_{\Delta,W}(Ph,Pk)
 \ge\frac{W}{8\Delta^2}
       \bigl(d_2(h,\mathcal F)-d_2(k,\mathcal F)\bigr)^2.
\end{equation}
\end{lemma}

\begin{proof}
For any nearest-point projection onto a closed convex subset of a
complete CAT(0) space, the projection is constant along the segment
from a point to its projection. To check this, let $q=Px$ and let
$z\in[q,x]$ satisfy $d(z,q)=\theta d(x,q)$. For $w$ in the convex set,
\[
 d(z,w)\ge d(x,w)-d(x,z)\ge\theta d(x,q)=d(z,q).
\]
Uniqueness gives $Pz=q$. Projection non-expansiveness and
\eqref{eq:fourpoint}, applied to the midpoints of $[x,Px]$ and
$[y,Py]$, now imply
\[
 d(x,y)^2-d(Px,Py)^2
   \ge\tfrac12\bigl(d(x,\mathcal F)-d(y,\mathcal F)\bigr)^2.
\]
The function $\psi(s)=\ell(s/\Delta)-s^2/(2\Delta^2)$ is convex,
since its second derivative on its finite domain is
$\Delta^{-2}\tan^2(s/\Delta)$. Decompose the kinetic cost as
\[
 \mathcal K_{\Delta,W}(h,k)
 =\frac{W}{4\Delta^2}d_2(h,k)^2
       +\frac W2\sum_i\psi(\sigma_i(h,k)).
\]
Use the preceding distance inequality for the quadratic term and
\eqref{eq:projectioncost} for the remainder.
\end{proof}

When applying these lemmas to a polygon, we project its nodes and then
join the projected nodes by geodesics. We do not require projection to
commute with geodesic interpolation. Fixed endpoint comparisons are
valid only when the endpoint norms belong to $\mathcal F$, unless the
changed endpoints are separately restored.

\subsection{Common bases and order retractions}
The next finite-dimensional observation permits spectral comparisons for
operations on a pair of norms whose apartments need not agree.

\begin{lemma}\label{lem:finitecommonbases}
Let $h(s)$ and $k(s)$, $0\le s\le1$, be split geodesics, each in its
own fixed apartment. There is a finite subdivision on each closed piece
of which both curves are diagonal in one fixed basis, with affine weights.
The same statement holds with real-analytic weights defined on a
neighborhood of $[0,1]$, with analytic weights in the common bases.
\end{lemma}

\begin{proof}
Write the two fixed bases as $e,f=eM$, with weights $a_i(s),c_j(s)$.
At a given parameter choose a nonzero pivot maximizing
\[
 \log|M_{ij}|+a_i(s)-c_j(s).
\]
Its column inequalities say
$|M_{kj}/M_{ij}|e^{a_k}\le e^{a_i}$.
Thus replacing $e_i$ by
$e_i+\sum_{k\ne i}(M_{kj}/M_{ij})e_k$ is an isometric triangular
change of the $h$-basis, and makes $f_j=M_{ij}e_i^{\rm new}$.
The row inequalities similarly give
$|M_{il}/M_{ij}|e^{c_j}\le e^{c_l}$; replacing
$f_l$ by $f_l-(M_{il}/M_{ij})f_j$ for $l\ne j$ is an isometric
change of the $k$-basis and eliminates the pivot row. The inverses
satisfy the same inequalities. Continue on the constant invertible
Schur complement.

There are finitely many complete pivot sequences. Each produces
constant basis matrices; its validity is a finite list of inequalities
in the original weights plus constants. For affine weights each
validity set is an interval, possibly empty or a point. Their union
covers $[0,1]$, and subdivision at the finitely many endpoints gives
the assertion. At completion the two final bases differ by a
permutation and a diagonal matrix, so the new weights differ from
the original weights only by fixed constants and permutation.

For analytic weights the differences defining validity are analytic
on a neighborhood of the closed interval. Each is either identically
zero or has finitely many zeros there. Subdivide at these zeros for
all the finitely many pivot sequences. Validity is constant on each
open piece and extends to its endpoints by non-strict inequalities.
This proves the analytic assertion as well. Ties require no separate
choice of a limiting basis.
\end{proof}

For two norms define their maximum and infimal convolution by
\[
 (h\vee k)(z)=\max\{h(z),k(z)\},\qquad
 (h\wedge k)(z)=\inf_{z=x+y}\max\{h(x),k(y)\}.
\]
The latter is a norm: equivalence to a fixed finite-dimensional norm
gives a positive lower bound on nonzero vectors, and approximate
minimizing decompositions give the ultrametric triangle inequality.
In a common orthogonal basis these operations are coordinatewise
maximum and minimum of logarithmic weights. Both are jointly
non-expansive for the maximum $d_\infty$ distance, as is clear from
multiplicative comparison of their defining formulas.

For an arbitrary pair $H=(h_A,h_C)$ define
\begin{equation}\label{eq:orderretraction}
 \Pi H=\bigl(\midpt(h_A,h_A\wedge h_C),
             \midpt(h_C,h_A\vee h_C)\bigr).
\end{equation}
For an ordered pair and $s\ge0$, also define
\begin{equation}\label{eq:shrinkandcap}
 \begin{aligned}
 \mathsf S_sH&=\bigl(\midpt(h_A,h_C\wedge e^s h_A),
                    \midpt(h_C,h_A\vee e^{-s}h_C)\bigr),\\
 \mathsf C_sH&=\bigl(\midpt(h_A,h_A\vee e^{-s}h_C),
                    \midpt(h_C,h_C\wedge e^s h_A)\bigr).
 \end{aligned}
\end{equation}
These formulas are intrinsic. In a common basis with weight pair
$(a_i,c_i)$, $\Pi$ leaves the pair unchanged when $a_i\le c_i$
and otherwise replaces it by its common average. For the two ordered
operations, put $m_i=(a_i+c_i)/2$ and $g_i=c_i-a_i\ge0$.
They retain $m_i$ and replace $g_i$, respectively, by
\begin{equation}\label{eq:gapoperations}
 (g_i-s)_+\quad\text{and}\quad\min\{g_i,s\}.
\end{equation}

\begin{lemma}\label{lem:pairretractions}
The map $\Pi$ retracts all norm pairs onto $\Dord$. The maps
$\mathsf S_s,\mathsf C_s$ preserve $\Dord$ and do not increase its
determinant gap. All three preserve $d(h_A)+d(h_C)$, preserve
splitness, and are continuous in the maximum $d_\infty$ metric.
For any one of these maps $T$, with $s$ fixed where applicable,
\begin{equation}\label{eq:totalcostcontraction}
 \sum_{v,i}\phi\bigl(\sigma_i((TH)_v,(TH')_v)\bigr)
 \le\sum_{v,i}\phi\bigl(\sigma_i(h_v,h'_v)\bigr)
\end{equation}
for every lower-semicontinuous convex $\phi$. The same scalar function
is used at both vertices.
\end{lemma}

\begin{proof}
The pointwise assertions follow from the common-basis formulas.
Joint $d_\infty$ continuity follows from the maximum, infimal
convolution and canonical midpoint formulas, and extends the assertions
to all norms by split approximation. Ordered pairs can be approximated
by split ordered pairs: after sufficiently close split approximation,
a vanishing opposite scalar shift restores order.

For the cost inequality first join $H$ to $H'$ by a product geodesic
on $[0,1]$. By Lemma~\ref{lem:finitecommonbases} it has finitely many
fixed common-basis pieces. Subdivide further where a signed gap is
zero, or where an ordered gap equals $s$. On every resulting piece
the transformed weights are affine and their derivatives either stay
unchanged or are replaced by their pairwise average. Jensen's
inequality gives
\[
 \sum_{v,i}\phi((Th)'_{vi})\le\sum_{v,i}\phi(h'_{vi}).
\]
The right-hand side is the constant endpoint spectral cost of the
original product geodesic. Apply the spectral triangle majorization
and Jensen's inequality to the transformed finite path. Its endpoint
cost is at most the integral of its derivative cost, proving
\eqref{eq:totalcostcontraction} for split inputs. For all-norm inputs
first pass to the limit for finite continuous convex $\phi$. A general
lower-semicontinuous convex function is the increasing supremum of
finite maxima of affine minorants; this gives the extended-valued
assertion without assuming continuity at a barrier.
\end{proof}

In particular the three maps contract the total log-cosine kinetic
cost between two nodes. The parameter $s$ in this statement is constant.
A time-dependent shrink introduces additional velocities, which will
be estimated separately in the proof of finite incoming contact.
